\begin{parts}

\part درست

 همگرایی سیاست تنها به \lr{argmax} مقادیر $Q$ در هر حالت بستگی دارد،
نه به مقدار دقیق آن‌ها. بنابراین ممکن است پس از چند ایتریشن، \lr{argmax} در همه‌ی حالت‌ها
ثابت شود و سیاست حریصانه دیگر تغییر نکند، در حالی که $V$ هنوز به مقدار دقیق $V^*$ نرسیده
و به‌روزرسانی آن ادامه دارد.

\part نادرست

 \lr{Q-learning} یک الگوریتم \lr{off-policy} است؛ یعنی مقادیر بهینه‌ی $Q$ را می‌تواند
مستقل از اینکه کنش‌های واقعاً انتخاب‌شده طبق سیاست رفتاری بهینه باشند یا نه، یاد بگیرد.
شرط کافی همگرایی، بازدید کافی از همه‌ی زوج‌های $(s,a)$ است، نه بهینه بودن سیاست رفتاری.

\part درست 

 در \lr{Value Iteration} هر ایتریشن شامل به‌روزرسانی همه‌ی حالت‌ها با پیچیدگی
$O(|S|^2|A|)$ است. در \lr{Policy Iteration} هر ایتریشن شامل ارزیابی سیاست با پیچیدگی
$O(|S|^3)$ و بهبود سیاست با پیچیدگی $O(|S|^2|A|)$ است.
اگر $|A|\approx |S|$، هر دو از مرتبه‌ی $O(|S|^3)$ می‌شوند، پس پیچیدگی هر ایتریشن تقریباً
برابر خواهد بود.

\part نادرست

 اگر به همه‌ی پاداش‌ها یک ثابت $c$ اضافه کنیم، سیاست بهینه لزوماً ثابت نمی‌ماند.
دلیلش این است که اگر \lr{MDP} حالت \lr{terminal} داشته باشد، سیاست‌هایی که دیرتر به
حالت پایانی می‌رسند تعداد پاداش‌های بیشتری می‌گیرند؛ پس مقدار اضافه‌شده برای همه‌ی سیاست‌ها
یکسان نیست.

البته اگر افق نامحدود باشد و پایان اپیزود در کار نباشد، در این حالت مقدار اضافه‌شده برای
همه‌ی سیاست‌ها برابر با
\[
\frac{c}{1-\gamma}
\]
می‌شود و آنگاه گزاره درست خواهد بود. اما در حالت کلی گزاره نادرست است.

\part نادرست

 مقدار $\gamma$ میزان اهمیت پاداش‌های آینده نسبت به پاداش‌های نزدیک را مشخص می‌کند.
پس اگر $\gamma$ عوض شود، ممکن است کنش بهینه هم عوض شود.

برای مثال یک \lr{MDP} را در نظر بگیرید که از حالت شروع $s$ دو کنش داریم:

\[
a_1:\quad \text{دریافت پاداش } 5 \text{ و پایان}
\]

\[
a_2:\quad \text{ابتدا پاداش } 0 \text{ و سپس در گام بعدی پاداش } 10 \text{ و پایان}
\]

ارزش کنش اول برابر است با:

\[
Q(s,a_1)=5
\]

و ارزش کنش دوم برابر است با:

\[
Q(s,a_2)=10\gamma
\]

اگر $\gamma=0.4$ باشد، داریم:

\[
Q(s,a_2)=4<5
\]

پس کنش $a_1$ بهتر است. اما اگر $\gamma=0.9$ باشد، داریم:

\[
Q(s,a_2)=9>5
\]

پس کنش $a_2$ بهتر می‌شود. بنابراین دو \lr{MDP} که فقط در مقدار $\gamma$ فرق دارند،
لزوماً سیاست بهینه‌ی یکسانی ندارند.

\part درست

 در افق محدود (\lr{finite horizon})، ارزش هر کنش به تعداد گام‌های باقی‌مانده تا پایان
افق بستگی دارد. بنابراین حتی برای یک حالت ثابت $s$، سیاست بهینه بسته به این‌که چند گام تا
پایان مانده می‌تواند فرق کند؛ یعنی سیاست بهینه \lr{non-stationary} است، برخلاف \lr{MDP}
با افق نامحدود که سیاست بهینه \lr{stationary} است.
\end{parts}


\vspace{1cm}

\begin{parts}

\part
تفاوت \lr{on-policy RL} و \lr{off-policy RL}:
در روش‌های \lr{on-policy}، عامل ارزش سیاستی را یاد می‌گیرد که برای تولید تجربه از همان
استفاده می‌کند؛ یعنی سیاست رفتاری (\lr{behavior policy}) و سیاست هدف (\lr{target policy})
یکسان هستند، مانند \lr{SARSA}. در روش‌های \lr{off-policy}، سیاست تولید تجربه می‌تواند با
سیاستی که عامل قصد یادگیری ارزش آن را دارد متفاوت باشد، مانند \lr{Q-learning} که $Q^*$ را
یاد می‌گیرد، حتی اگر کنش‌ها با سیاستی مثل \lr{$\varepsilon$-greedy} انتخاب شده باشند.

\part
تفاوت \lr{model-based RL} و \lr{model-free RL}:
در \lr{model-based RL}، عامل مدل محیط یعنی توابع گذار $T(s'\mid s,a)$ و پاداش
$R(s,a,s')$ را از قبل می‌داند یا از تجربه تخمین می‌زند و بعد با استفاده از این مدل
برنامه‌ریزی می‌کند. در \lr{model-free RL}، عامل بدون ساختن مدل صریحی از محیط، مستقیماً
از تعامل با محیط تابع ارزش یا سیاست را یاد می‌گیرد، مانند \lr{Q-learning}، \lr{SARSA}
و \lr{Monte Carlo}.

\end{parts}